\chapter{任务不交还处理器时怎么办}

\begin{keyidea}
常驻协调者已经能保存、选择和恢复 A/B，每项任务也有独立现场记录与栈；但切换入口只能
由正在运行的任务主动调用。要限制一项任务占用处理器的时间，机器必须加入与当前指令流
独立的事件，并在事件入口形成可恢复现场。
\end{keyidea}

\section{常驻并不等于有机会执行}

B 执行下面的指令序列：

\begin{lstlisting}
repeat_forever:
  compute
  jump repeat_forever
\end{lstlisting}

处理器会不断从 B 的地址取指。协调代码虽然存在于另一段内存，但“存在”不会自动改变指令
位置。外部观察者也不能靠在内存里写一个标志解决问题，因为 B 没有读取该标志。

要在 B 不执行交还指令时重新获得控制，机器必须提供一种与当前指令流独立的事件：事件到达
后，处理器先保存足以继续 B 的状态，再把指令位置改成预先登记的入口。

\section{为什么选择可编程时钟事件}

设备完成也能产生外部事件，但机器空闲时未必有设备活动，不能用它保证处理器一定被收回。
一个可编程计时器可以在设定时间后发出事件，与 B 是否读内存、是否调用函数无关。

\subsection{新器件接到现有机器的哪里}

第一章的时钟源只能推动电路状态变化，软件不能要求它在未来某个计数值通知处理器。现在
增加的计时器是一项独立器件：它接收同一参考时钟，通过地址互连暴露配置寄存器，并用一根
事件请求线连接处理器的异步事件输入。

本样机只有一个事件源，所以计时器请求线直接连接处理器，不增加中断控制器。以后事件源
增多，才会出现“多根请求线如何编号、屏蔽和路由”的新问题。

\begin{center}
\begin{tikzpicture}[
  node distance=11mm and 18mm,
  hw/.style={stage,minimum width=3.8cm},
  irq/.style={->,>=Latex,thick,draw=BookNavy},
  request/.style={->,>=Latex,thick,draw=BookAccent},
  response/.style={->,>=Latex,thick,draw=BookMuted}
]
\node[hw] (cpu) {扩展后的处理器\\异步事件输入};
\node[hw,right=of cpu] (bus) {原有地址互连\\增加计时器范围};
\node[hw,right=of bus] (timer) {可编程计时器\\计数、比较、pending};
\node[hw,above=13mm of timer] (clock) {原有参考时钟\\周期边沿};
\draw[request] ([yshift=3mm]cpu.east) -- node[above,font=\footnotesize]{寄存器读写请求}
  ([yshift=3mm]bus.west);
\draw[response] ([yshift=-3mm]bus.west) -- node[below,font=\footnotesize]{寄存器响应}
  ([yshift=-3mm]cpu.east);
\draw[request] ([yshift=3mm]bus.east) -- node[above,font=\footnotesize]{译码后的请求}
  ([yshift=3mm]timer.west);
\draw[response] ([yshift=-3mm]timer.west) -- node[below,font=\footnotesize]{计时器响应}
  ([yshift=-3mm]bus.east);
\draw[irq] (timer.south) -- ++(0,-11mm) -|
  node[below,font=\footnotesize,pos=0.5]{event request} (cpu.south);
\draw[flow] (clock) -- (timer);
\end{tikzpicture}
\end{center}
\diagramnote{总线方向用于配置和查询，事件请求线用于计时器主动通知处理器。软件写过寄存器
并不等于事件已经发生；只有计数条件满足后，计时器才置 pending 并拉起请求线。}

\subsection{计时器的硬件模型}

计时器内部保存四类最小状态：

\begin{longtable}{|L{0.22\linewidth}|L{0.28\linewidth}|L{0.38\linewidth}|}
\hline
\rowcolor{BookTableHead}
内部状态 & 状态变化 & 对外作用 \\ \hline
\code{counter} & 每个参考时钟边沿递增 & 提供单调增加的比较基准 \\ \hline
\code{deadline} & 处理器通过写事务设置 & 指定下一次到期计数值 \\ \hline
\code{enabled} & 处理器启用或停止计时比较 & 决定是否检查到期条件 \\ \hline
\code{pending} & 到期时置 1，确认后清 0 & 为 1 时保持事件请求有效 \\ \hline
\end{longtable}

这四类状态通过比较器和置位/清除逻辑连接，而不是由软件循环反复读取后自行判断。参考时钟
只递增 \code{counter}；比较器同时观察 \code{counter}、\code{deadline} 与
\code{enabled}；比较成立时，它向 \code{pending} 发出置位条件。软件确认走独立清除
输入，因此一次很短的“到期相等”可以转换为持续有效的事件请求。

\begin{center}
\begin{tikzpicture}[node distance=9mm and 6mm]
\node[stage,minimum width=2.35cm] (clock) {参考时钟\\周期边沿};
\node[stage,minimum width=2.55cm,right=of clock] (counter) {\code{counter}\\每边沿递增};
\node[stage,minimum width=2.75cm,right=of counter] (compare) {比较与使能门\\\(counter\ge deadline\)};
\node[stage,minimum width=2.55cm,below=13mm of compare] (pending) {\code{pending}\\置位后保持};
\node[stage,minimum width=2.5cm,right=of pending] (event) {事件请求线\\送往处理器};
\node[stage,minimum width=2.6cm,above=11mm of compare] (deadline)
  {\code{deadline}\\软件写入的比较值};
\node[stage,minimum width=2.6cm,below=11mm of counter] (control)
  {\code{enabled}\\软件写入的使能位};
\node[stage,minimum width=2.6cm,below=11mm of pending] (ack)
  {\code{STATUS/ACK}\\写 1 产生清除条件};
\draw[flow] (clock) -- (counter);
\draw[flow] (counter) -- (compare);
\draw[flow] (deadline) -- (compare);
\draw[flow] (control) -- (compare);
\draw[flow] (compare) -- node[above,font=\scriptsize]{到期时置位} (pending);
\draw[->,>=Latex,thick,BookMuted] (ack) -- node[right,font=\scriptsize]{确认后清除} (pending);
\draw[flow] (pending) -- (event);
\end{tikzpicture}
\end{center}
\diagramnote{比较器只产生置位条件，\code{pending} 才保存尚未处理的事件。请求线直接反映
\code{pending}，所以处理器暂时未进入事件入口时，通知不会因比较条件只成立一个边沿而丢失。}

每个时钟边沿后，若 \code{enabled=1} 且 \code{counter >= deadline}，硬件就设置
\code{pending=1}。请求线保持有效，直到软件确认事件。采用“保持”而不是一个极短脉冲，
可以避免处理器暂时关闭事件接收时永久丢失通知。

这个模型没有自动周期。协调代码处理一次到期事件后，写入新的 deadline，才形成重复时间
片。这样可以先看清“本次到期”和“安排下次到期”是两个动作。

\subsection{计时器的软件模型}

软件不直接访问内部连线，而是通过地址互连看到一组内存映射寄存器：

\begin{longtable}{|L{0.33\linewidth}|L{0.17\linewidth}|L{0.38\linewidth}|}
\hline
\rowcolor{BookTableHead}
物理地址与名称 & 访问方式 & 软件可见语义 \\ \hline
\code{0x40001000 COUNT} & 只读 & 返回当前 counter 快照 \\ \hline
\code{0x40001008 DEADLINE} & 读写 & 设置下一次到期比较值 \\ \hline
\code{0x40001010 CONTROL} & 读写 & bit 0 控制 enabled \\ \hline
\code{0x40001014 STATUS/ACK} & 读/写 & 读 bit 0 得 pending；写 1 清除 pending \\ \hline
\end{longtable}

\code{0x40001000} 恰好位于 UART 保留窗口之后，等于
\(\code{0x40000000}+\code{0x1000}\)。这里的 \code{0x1000} 是一页 4 KiB，
所以两个设备拥有互不重叠、分别页对齐的译码窗口；以后即使增加按页权限、缓存属性或
总线跟踪，也不必把两个设备拆出同一页。这仍是教学样机的选择，不是计时器标准常量。
换用别的空闲窗口没有原理障碍，但硬件译码与驱动地址必须作为一份接口契约同步更新。

\code{COUNT} 与 \code{DEADLINE} 都是 64 位，因而分别从 \(+0\) 和 \(+8\)
开始；两个 32 位控制寄存器接着位于 \(+0x10\) 与 \(+0x14\)。最后一个偏移采用
“读状态、写 1 确认”的双重语义，是为了让观察 pending 与清除同一 pending 对准一个
控制点，并不表示它是普通可读写 RAM。读操作取得控制器状态快照，写操作改变控制器
状态；两者都没有把数据块在 RAM 与计时器之间来回搬运。

让计时器在 \(Q\) 个计数单位后产生一次事件，软件按以下顺序配置：

\begin{lstlisting}
now = read64(COUNT)
write64(DEADLINE, now + Q)
write32(STATUS_ACK, 1)       // clear an earlier pending event
write32(CONTROL, ENABLE)
\end{lstlisting}

若先启用、后写 deadline，旧 deadline 可能立即满足并产生一次非预期事件。寄存器顺序因此
是软件模型的一部分，不只是代码风格。

\subsection{计时器提供的抽象}

上层软件得到“在单调计数达到 deadline 后形成一个保持到确认的事件”。它没有得到以下
保证：

\begin{itemize}
\item deadline 不是日期，也不自动等于毫秒；
\item 到期表示请求已经提出，不表示处理器恰好在该计数值执行入口；
\item 事件被屏蔽时，请求可以延迟交付；
\item 计时器不会保存任务现场，也不会选择下一项任务。
\end{itemize}

因此，计时器抽象适合保证协调代码最终重新获得一次入口机会；任务选择仍属于软件。

\subsection{处理器异步入口也需要明确三层模型}

计时器有了请求线，原处理器仍只会顺序取指。机器还要扩展处理器的事件入口能力。

\textbf{硬件模型。} 处理器在一条指令完成后检查事件输入。若事件允许且请求有效，处理器
保存下一条指令位置、原栈指针和状态标志，切到固定入口栈顶，关闭同类事件嵌套，再把指令
位置改为固定事件入口。事件返回指令执行相反的受控恢复。

\textbf{软件模型。} 固件把事件入口约定为 \code{0x00008000}，把入口栈顶约定为
\code{0x001f0000}。硬件建立的最小帧按顺序包含：

\begin{lstlisting}
event_frame:
  resume_instruction
  interrupted_stack_pointer
  saved_flags
\end{lstlisting}

入口代码能读取这份帧，并通过确认寄存器清除计时器 pending。普通通用寄存器尚未由硬件
保存，必须由入口软件补齐。

\textbf{抽象。} 软件得到一个可在普通指令流之外发生的受控入口，以及一份足以返回被打断
位置的最小继续状态。该抽象不等于完整上下文切换：它不保存全部寄存器，也不决定事件到达
后继续原任务还是恢复另一个任务。

处理器扩展后的硬件路径可以分成五个相邻结构。请求锁存器记住计时器通知；边界检查器只在
一条普通指令已经提交且事件允许时接受它；最小帧写入器把继续位置、旧 SP 和标志写到入口
栈；入口选择器再把 PC 与 SP 改为固定入口值。事件返回译码器只接受符合约定的帧，并按
相反方向恢复这些字段。任何一步失败都不能假装已经进入协调代码。

\begin{center}
\begin{tikzpicture}[node distance=9mm and 3mm]
\node[stage,minimum width=2.3cm] (request) {事件请求锁存\\保持待处理来源};
\node[stage,minimum width=2.45cm,right=of request] (boundary) {指令边界检查\\事件是否允许进入};
\node[stage,minimum width=2.5cm,right=of boundary] (frame) {最小帧写入器\\继续 PC、旧 SP\\与标志};
\node[stage,minimum width=2.4cm,below=12mm of frame] (select) {入口 PC/SP 选择\\切到固定入口};
\node[stage,minimum width=2.3cm,right=of select] (fetch) {入口取指\\协调代码开始};
\draw[flow] (request) -- (boundary);
\draw[flow] (boundary) -- (frame);
\draw[flow] (frame) -- (select);
\draw[flow] (select) -- (fetch);
\node[stage,fill=BookWarm,minimum width=6.8cm,below=12mm of select] (return)
  {事件返回译码器：检查帧格式\\恢复 PC、SP 与标志，再重新开放事件};
\draw[flow] (fetch.south) |- (return.east);
\end{tikzpicture}
\end{center}
\diagramnote{异步入口不是把请求线直接接到 PC。处理器先等待精确指令边界，再建立可返回的
最小帧，最后切换入口 PC/SP；事件返回只有在帧检查通过后才恢复被选中的执行流。}

为支持这条路径，硬件至少要明确：

\begin{longtable}{|L{0.21\linewidth}|L{0.30\linewidth}|L{0.37\linewidth}|}
\hline
\rowcolor{BookTableHead}
硬件能力 & 必须提供的事实 & 软件随后才能做的事 \\ \hline
计时器 & 到期时产生一个可识别事件 & 规定最长连续运行区间 \\ \hline
事件入口表 & 事件编号对应固定入口地址 & 让协调代码得到控制 \\ \hline
最小现场保存 & 保存被打断继续位置与必要标志 & 以后从原指令流继续 \\ \hline
事件屏蔽状态 & 防止入口最早阶段被同类事件反复覆盖 & 完成一次一致的现场建立 \\ \hline
事件返回指令 & 从规定帧恢复状态与执行级别 & 回到所选任务 \\ \hline
\end{longtable}

硬件不必替软件决定运行 A 还是 B。它只保证事件发生时能离开 B，并把控制交到一个软件已
知入口。选择仍由协调代码完成。

\subsection{硬件保存的状态为什么通常不够}

事件到达的第一目标是安全进入入口，而不是一次完成所有任务切换。处理器通常只自动保存
继续位置、状态标志以及进入和返回所必需的少量字段。通用寄存器若要保持不变，入口软件仍
要尽早保存。

\begin{center}
\begin{tikzpicture}[node distance=10mm and 15mm]
\node[stage] (b) {B 的任意指令\\计时器到期};
\node[stage, right=of b] (hw) {硬件建立最小帧\\改到事件入口};
\node[stage, right=of hw] (sw) {入口补存寄存器\\形成 B 的完整现场};
\node[stage, below=13mm of sw] (choose) {协调者判断\\A 是否可运行};
\node[stage, left=of choose] (restore) {恢复 A 的现场\\设置下次计时};
\node[stage, left=of restore] (a) {事件返回\\A 继续执行};
\draw[flow] (b) -- (hw);
\draw[flow] (hw) -- (sw);
\draw[flow] (sw) -- (choose);
\draw[flow] (choose) -- (restore);
\draw[flow] (restore) -- (a);
\end{tikzpicture}
\end{center}
\diagramnote{时钟事件只强制打开入口机会。完整保存、任务选择和恢复仍是软件协议。}

入口不能先使用尚未保存的通用寄存器计算，因为那会覆盖 B 的值。最早几条指令只能使用
硬件约定允许破坏的状态，或者把寄存器按固定顺序压入已知安全的栈。

\subsection{事件入口使用哪一份栈}

B 的栈可能已接近边界，甚至因为错误把栈指针改成无效值。若硬件入口无条件继续使用 B 的
栈，保存第一项状态就可能失败。机器需要一个不由普通任务控制的入口栈位置；单核样机可让
硬件或最早入口从固定协调状态取得它。

\begin{longtable}{|L{0.22\linewidth}|L{0.30\linewidth}|L{0.36\linewidth}|}
\hline
\rowcolor{BookTableHead}
栈 & 保存的调用关系 & 所有权边界 \\ \hline
A 的任务栈 & A 的函数返回地址与局部数据 & A 运行或暂停时保留 \\ \hline
B 的任务栈 & B 的函数返回地址与局部数据 & B 运行或暂停时保留 \\ \hline
协调入口栈 & 事件最早保存与选择代码局部状态 & 普通任务不应改写 \\ \hline
\end{longtable}

入口栈解决的是“如何安全开始保存”，不是自动隔离。由于当前机器仍允许 B 写任意物理地址，
B 依然能够故意或错误地覆盖协调入口栈；这个缺陷留到章末。

\subsection{怎样把事件帧并入 B 的任务记录}

事件帧记录 B 被打断的位置，软件补存通用寄存器。协调者可以让 B 的任务记录指向这份完整
帧，而不必把每个字段再复制一次。关键是记录必须说明帧位于哪一段栈、布局采用哪个入口
约定。

\begin{lstlisting}
timer_entry:
  hardware has saved resume position and flags
  save remaining registers in fixed order
  current_task.saved_stack = current stack pointer
  current_task.status = ready
  switch to coordinator stack
  next = choose_ready_task()
  restore next.saved_stack
  return_from_event
\end{lstlisting}

主动交还与时钟打断可以最终汇入相同的“已形成完整现场”边界，但两者最早入口不同。主动
交还遵守普通调用约定；时钟可以发生在任意指令，必须按异步入口约定保存所有可能被后续代码
破坏的状态。

\section{第一次出现可运行集合}

只有两项任务时也需要区分“暂停但能继续”和“已经完成”。协调者不应把完成的任务重新装入
处理器。可以维护一个很小的可运行集合：

\begin{longtable}{|L{0.18\linewidth}|L{0.28\linewidth}|L{0.42\linewidth}|}
\hline
\rowcolor{BookTableHead}
任务状态 & 是否参与选择 & 状态改变的入口 \\ \hline
ready & 是 & 被时钟打断或主动交还后 \\ \hline
running & 否，当前已占用处理器 & 协调者恢复该任务时 \\ \hline
finished & 否 & 任务进入完成入口时 \\ \hline
not started & 可先构造初始现场再加入 & 装入者完成代码、数据和栈准备时 \\ \hline
\end{longtable}

选择规则可以是 A、B 轮流。每次恢复前重新设置计时器，给所选任务一个有限运行区间。时间
区间结束并不表示任务完成，只表示协调者获得一次重新选择的机会。

\subsection{“被唤醒”和“正在运行”尚未出现}

本章任务只做计算，没有等待设备或条件；ready 已足够表达“具有运行资格”。以后加入等待
后，某项任务即使没有完成也可能暂时不应进入选择集合，届时才需要额外状态和唤醒协议。
现在提前画完整生命周期只会让尚无问题来源的概念进入样机。

\subsection{一次强制切换的逐时刻状态}

\begin{longtable}{|L{0.11\linewidth}|L{0.27\linewidth}|L{0.25\linewidth}|L{0.25\linewidth}|}
\hline
\rowcolor{BookTableHead}
时刻 & 控制流位置 & B 的最新现场 & A 的最新现场 \\ \hline
\(t_0\) & B 的循环 & 处理器中 & A 的内存记录与栈 \\ \hline
\(t_1\) & 硬件事件入口 & 最小字段在事件帧，其余仍在寄存器 & 不变 \\ \hline
\(t_2\) & 软件入口完成保存 & B 的完整帧与栈 & 不变 \\ \hline
\(t_3\) & 协调者选择 A & 稳定，可供以后恢复 & 正在被读取 \\ \hline
\(t_4\) & 事件返回到 A & B 的内存现场 & 处理器中 \\ \hline
\end{longtable}

事件返回的目标不必是刚才被打断的 B。硬件要求返回帧格式正确；软件可以在返回前换成 A
的已保存帧。这个能力使时钟入口成为抢占切换点。

\section{本章新增结构的准确边界}

这一阶段只新增：

\begin{itemize}
\item 可编程计时器及其事件入口；
\item 不依赖普通任务栈的最早入口状态；
\item 能容纳异步打断的完整任务现场格式；
\item ready/running/finished 三类选择状态；
\item 每次恢复前重新设置时间预算的轮流选择规则。
\end{itemize}

它没有新增文件、用户态、系统调用、虚拟地址、权限或设备服务。所谓“最小抢占调度”在此
只表示：协调者能在任务未主动交还时获得入口，并从可运行集合选择另一项任务。

\topic{下一项问题}
B 已经不能永久独占处理器，却仍能执行一条写指令覆盖 A 的栈、
A 的代码、协调者记录或时钟入口。时间上的轮流使用已经建立，空间上的互不破坏仍未建立。
下一次演进将从这次真实覆盖失败出发，建立地址转换、访问权限和受保护的控制入口。
